\begin{center}\begin{tikzpicture}[font=\small,>=stealth,line join=round]
\draw[gray!40,rounded corners=3pt] (0,0) rectangle (14.2,9.0);
\node at (7.1,8.4) {例5：上下离中线同为17，未知函数值随即确定};
\begin{scope}[shift={(3.8,5.6)},x=.7cm,y=.16cm]\draw[->,gray] (-4,0)--(4.1,0) node[right]{$x$};\draw[->,gray] (0,-27)--(0,13) node[left]{$y$};\draw[teal,dashed] (-4,-7)--(4,-7);\fill[blue] (3,10) circle(2pt) node[above left]{$(3,10)$};\fill[orange!85!black] (-3,-24) circle(2pt) node[below right]{$(-3,-24)$};\draw[gray,dotted] (3,10)--(-3,-24);\fill[teal] (0,-7) circle(2pt) node[below right]{$(0,-7)$};\draw[<->,blue] (3.6,-7)--(3.6,10) node[midway,right]{$17$};\draw[<->,orange!85!black] (-3.6,-24)--(-3.6,-7) node[midway,left]{$17$};\end{scope}\node[align=left,anchor=west] at (7.8,4.5) {$g(x)=f(x)+7$ 为奇函数\\[12pt]$g(3)=17,\ g(-3)=-17$\\[12pt]$f(3)+f(-3)=-14$\\[12pt]$f(-3)=-24$};
\node[text width=12.8cm,align=center] at (7.1,.65) {只画已知与必然的点及辅助线，不假设参数 $a,b$ 决定的具体曲线。};
\end{tikzpicture}\end{center}
